\section{The determinant orbit interface, with its exact obstruction scope}
\label{sec:gct-obstruction-interface}

All varieties and representations in this section are over $\mathbb C$.
Put $V_D=\operatorname{Mat}_{D\times D}(\mathbb C)$,
$W_D=\operatorname{Sym}^D(V_D^*)$, and $G_D=\operatorname{GL}(V_D)$.
The action, including its inverse, is
\[
 (g\cdot f)(X)=f(g^{-1}X),\qquad
 \Omega_D=\overline{G_D\cdot\det_D}\subset W_D.
\]
Closures below are Zariski closures. The affine coordinate ring of $W_D$
is $B_D=\operatorname{Sym}(W_D^*)$, with degree-$e$ part
$\operatorname{Sym}^e(\operatorname{Sym}^D V_D)$.
Here $D$ is the matrix size; $e$ is a coordinate-ring degree and is
not the polynomial degree $D$.

\subsection{Original padding and stabilizer conventions}

GCT I, Section 4, Proposition 4.1, constructs the determinant orbit-closure
map from a formula through a possibly singular linear substitution;
GCT II, Section 3.1, uses an independent padding variable
\cite{MS2001Interface,MS2006Interface}. Explicitly, for $m<D$, let
$X$ be the bottom-right $m\times m$ submatrix of $Y\in V_D$ and
$y=Y_{11}$, which is outside that submatrix. The linear map on forms is
\[
 j_{D,m}:\operatorname{Sym}^m(V_m^*)\longrightarrow W_D,
 \qquad h\longmapsto y^{D-m}h(X).
\]
It is injective: the ring substitution retaining $X$ and putting $y=1$
and all other entries of $Y$ equal to zero is a left inverse on its
image. There is also a homogeneous coefficient left inverse, obtained
by taking the coefficient of $y^{D-m}$ and then setting the other
outside variables to zero. Thus this map has no kernel and explicitly
retains every coefficient of $h$. In the notation of GCT II, its $n$
is the present $m$, and its $m$ is the present $D$.

Those papers use projective $\operatorname{SL}(V_D)$ orbit closures.
The precise comparison with the affine convention above is as follows.
For $g\in G_D$, choose $\mu\in\mathbb C^*$ with
$\mu^{D^2}=\det(g)$ and set $s=\mu^{-1}g$. Then
$s\in\operatorname{SL}(V_D)$ and
$g\cdot f=\mu^{-D}(s\cdot f)$. Hence the projective orbits agree.
Moreover every nonzero scalar multiple of $f$ is in $G_D\cdot f$,
because $\lambda I$ acts by $\lambda^{-D}$. The affine closure is
therefore a cone, and its homogeneous vanishing ideal is precisely the
homogeneous ideal of the projective orbit closure: a homogeneous equation
vanishes on the projective orbit exactly when it vanishes on every
nonzero vector above it. This proves the comparison of coordinate rings,
with their degrees retained.

The original stabilizer discussion distinguishes a form from its line.
For example the linear substitution $T_{A,B}(Y)=AYB^{-1}$, with
$A,B\in\operatorname{GL}_D$, satisfies
\[
 \det(T_{A,B}Y)=\frac{\det A}{\det B}\det Y,
 \qquad
 T_{A,B}\cdot\det_D=\frac{\det B}{\det A}\det_D.
\]
The first identity follows by multiplication of determinants and the
second by inserting $T_{A,B}^{-1}$. Thus every such substitution fixes
the determinant line, and it fixes the form exactly when
$\det A=\det B$. Transposition fixes the form as well. To intersect
these substitutions with the original ambient $\operatorname{SL}(V_D)$,
retain the additional equation
$(\det A/\det B)^D=1$; with transposition the left side is multiplied
by $(-1)^{D(D-1)/2}$. Indeed left multiplication acts independently on
$D$ columns, right multiplication on $D$ rows, and transposition exchanges
the $D(D-1)/2$ off-diagonal coordinate pairs. These equations prevent
identifying the line stabilizer with the form stabilizer.
For the permanent, row and column permutations and transposition preserve
the form, and row scalars $\alpha_i$ and column scalars $\beta_j$
multiply it by $\prod_i\alpha_i\prod_j\beta_j$: each defining monomial
uses each row and column once. These verified subgroups supply the
stabilizer context used here; no classification of a new fibre stabilizer
is inferred from them.

\subsection{An explicit endomorphism and a strict boundary point}

For the retained-fibre family take $n\ge1$, $D=4n+1$, and
$E_n=\mathbb C^{5n+1}$ with ordered coordinates
\[
 (a_1,b_1,c_1,r_1,t_1,\ldots,a_n,b_n,c_n,r_n,t_n,z).
\]
Write $L_n:E_n\to V_D$ for the linear map whose matrix value is
the previously constructed $\mathcal M_n^h$, retaining all entries:
\[
 T_i^h=\begin{pmatrix}z&-r_i&0&0\\0&z&-r_i&0\\0&0&z&-r_i\\0&0&0&z\end{pmatrix},
 \quad v_i^h=\begin{pmatrix}-2a_i\\b_i\\-2z\\c_i\end{pmatrix},
 \quad
 \mathcal M_n^h=\begin{pmatrix}\operatorname{diag}(T_i^h)&v^h\\-u&0\end{pmatrix},
\]
where $u=(t_1,0,0,0,\ldots,t_n,0,0,0)$ and $v^h$ stacks the
$v_i^h$. Its determinant, already proved by block multiplication, is
\[
 H_n=z^{4n-4}\sum_{i=1}^n t_i
 (c_ir_i^3-2z^2r_i^2+b_ir_iz^2-2a_iz^3).
\]
It has degree $D$. Every coordinate of $E_n$ is a nonzero scalar
multiple of an entry of $L_n$, so $L_n$ is injective.

Order the $D^2$ entries of $Y\in V_D$ row by row. Let
$p_n:V_D\to E_n$ take the first $5n+1$ entries in exactly the displayed
order, and let $i_n:E_n\to V_D$ fill those entries and put zero in
the remaining positions. The inequality $5n+1<D^2$ holds for $n\ge1$,
and $p_ni_n=I$. These are specified linear maps, rather than an
identification of two variable spaces. Set
\[
 A_n=L_np_n\in\operatorname{End}(V_D),\qquad
 \widetilde H_n=H_n\circ p_n=\det_D\circ A_n\in W_D.
\]
On coordinate rings, $p_n^*$ is injective, $i_n^*p_n^*=I$, and
$L_n^*(X_{pq})=(\mathcal M_n^h)_{pq}$. The latter is the previously
proved surjection $\vartheta_n$, with its stated ideal of linear
relations as kernel. Thus $A_n^*=p_n^*L_n^*$ retains that kernel,
and $i_n^*(\widetilde H_n)=H_n$ recovers the full determinant polynomial.

\begin{proposition}\label{prop:retained-determinant-boundary}
For every $n\ge1$, $\widetilde H_n\in\Omega_{4n+1}$, and
$\widetilde H_n\notin G_{4n+1}\cdot\det_{4n+1}$.
\end{proposition}
\begin{proof}
Let $A_n(\epsilon)=A_n+\epsilon I_{V_D}$. Its determinant as a
$D^2\times D^2$ endomorphism is a monic polynomial of degree $D^2$
in $\epsilon$. Away from its finite root set it is invertible, and
\[
 f_\epsilon(Y)=\det_D(A_n(\epsilon)Y)
             =A_n(\epsilon)^{-1}\cdot\det_D(Y)
\]
belongs to the determinant orbit. The coefficients of $f_\epsilon$
are polynomials in $\epsilon$. Every equation vanishing on that orbit,
when evaluated on $f_\epsilon$, is consequently a polynomial vanishing
for all but finitely many complex values of $\epsilon$. It is zero
identically and hence at $\epsilon=0$. This proves the claimed closure
membership, without assigning an inverse to $A_n(0)$.

For strictness, define the linear subspace
$\operatorname{Ann}(f)=\{v:\partial_v f=0\}$, using formal directional
derivatives. The chain rule gives
$\operatorname{Ann}(g\cdot f)=g\operatorname{Ann}(f)$.
For the universal determinant,
$\partial_v\det_D=\sum_{i,j}v_{ij}\operatorname{cof}_{ij}(Y)$.
The cofactors are linearly independent: each is nonzero and has row
multidegree one except in row $i$, and column multidegree one except
in column $j$. Distinct pairs $(i,j)$ have distinct combined
multidegrees, so their monomials cannot cancel. Therefore
$\operatorname{Ann}(\det_D)=0$. In contrast,
$\ker p_n\subseteq\operatorname{Ann}(\widetilde H_n)$ and
$\dim\ker p_n=D^2-(5n+1)>0$. The orbit membership is impossible.
Finally $H_n$ is nonzero: set $z=t_1=c_1=r_1=1$, $a_1=b_1=0$
and $t_i=0$ for $i>1$, obtaining $H_n=-1$.
\end{proof}

\subsection{The coordinate quotient and the obstruction direction}

Put $Z_n=\overline{G_D\cdot\widetilde H_n}$. The proposition and
$G_D$-invariance of $\Omega_D$ give $Z_n\subseteq\Omega_D$.
Both are cones by the scalar calculation above. For any such closed
cone $Z$, its ideal $I(Z)\subset B_D$ is homogeneous: if
$F=\sum_eF_e$ vanishes on $Z$, then for $z\in Z$ the polynomial
$\sum_e t^eF_e(z)$ vanishes for every $t\in\mathbb C$, so every
$F_e(z)=0$. Consequently restriction is the explicit graded map
\[
 \rho_n:\mathbb C[\Omega_D]=B_D/I(\Omega_D)
 \longrightarrow B_D/I(Z_n)=\mathbb C[Z_n],\qquad
 F+I(\Omega_D)\longmapsto F+I(Z_n).
\]
It is well defined since $I(\Omega_D)\subseteq I(Z_n)$, is surjective
in each degree by taking a polynomial representative, and has kernel
$I(Z_n)/I(\Omega_D)$. For the action on functions
$(g\cdot F)(f)=F(g^{-1}\cdot f)$, stability of both sets proves
$\rho_n(g\cdot F)=g\cdot\rho_n(F)$ by evaluation at every point.

In each degree the quotient also gives the exact multiplicity inequality
\[
 \dim\operatorname{Hom}_{G_D}(S,\mathbb C[Z_n]_e)
 \le \dim\operatorname{Hom}_{G_D}(S,\mathbb C[\Omega_D]_e)
\]
for every irreducible finite-dimensional rational $G_D$-module $S$.
Here is the splitting argument. Choose a linear section of $(\rho_n)_e$
and average $u^{-1}s u$ over the compact unitary group $U(D^2)$ with
its invariant probability measure. The averaged map is a section since
$\rho_n$ is equivariant, and is unitary-equivariant by change of integration
variable. It is $G_D$-equivariant: differentiating its intertwining
identity on skew-Hermitian matrices and taking complex linear spans
gives the identity on all matrices, and exponentiating gives it on
$\operatorname{GL}_{D^2}(\mathbb C)$. Every invertible complex matrix has a logarithm,
as follows block by block from its Jordan form and the finite logarithm
series on the nilpotent part. Composing maps $S\to\mathbb C[Z_n]_e$
with this section is an injection into the other Hom space, proving
the inequality. This proof also shows exactly why a multiplicity in the
smaller orbit closure exceeding the corresponding determinant multiplicity
would obstruct inclusion in a different instance.

The stabilizer restriction has its own exact map. For $f\in W_D$ let
$H_f=\{h\in G_D:h\cdot f=f\}$. The orbit map
$\alpha_f(g)=g\cdot f$ induces
\[
 \alpha_f^*:\mathbb C[\overline{G_D\cdot f}]
 \lhook\joinrel\longrightarrow \mathbb C[G_D]^{H_f,\mathrm{right}},
 \qquad F\longmapsto(g\mapsto F(g\cdot f)).
\]
Regularity follows from the entries of $g^{-1}$ being regular on $G_D$;
right invariance follows from $gh\cdot f=g\cdot f$; injectivity follows
because a function vanishing on the orbit vanishes on its closure.
For $e$-homogeneous $F$ and a line stabilizer with
$h\cdot f=\chi_f(h)f$, the precise identity instead is
$\alpha_f^*F(gh)=\chi_f(h)^e\alpha_f^*F(g)$.
Thus form stabilizers give invariants and line stabilizers give this
specified character. No surjectivity onto all right-invariant functions
is asserted; that would require extension across the orbit boundary.

\subsection{The later no-occurrence result at its stated scope}

For a separate pair of matrix sizes $N>m$, B\"urgisser--Ikenmeyer--Panova
use the top-left projection $\pi:V_N\to V_m$ and the form
\[
 p_{N,m}(Y)=Y_{11}^{N-m}\operatorname{per}_m(\pi Y),\qquad
 Z^{\mathrm{BIP}}_{N,m}=\overline{\operatorname{GL}(V_N)\cdot p_{N,m}}.
\]
Here $Y_{11}$ is one of the permanent variables. Their source theorem,
\cite[Theorem 1.4]{BIP2018Interface}, states that for positive integers
$N,e,m$ with $N\ge m^{25}$ and $\lambda\vdash Ne$, occurrence of
$\lambda$ in $\mathbb C[Z^{\mathrm{BIP}}_{N,m}]$ implies occurrence
in $\mathbb C[\Omega_N]$. This is their proved theorem, not a result
derived from the present fibres. We retain their same-variable padding
and exponent $25$; the independent padding above is not silently
substituted into that statement. The central scalar $tI$ acts on
$\mathbb C[W_N]_e$ by $t^{Ne}$, explaining the partition size $Ne$.
The theorem compares occurrence, and does not assert the multiplicity
inequality needed to force a quotient.

A later concrete multiplicity example has a different, specified ambient
space. In $\mathbb C[x_1,x_2,x_3]_6$ put
$\mathrm{Ch}_3^6=\{\ell_1\cdots\ell_6\}$ and
$\mathrm{Pow}_{3,k}^6=\overline{\{\ell_1^6+\cdots+\ell_k^6\}}$,
where all $\ell_i$ are homogeneous linear forms. D\"orfler--Ikenmeyer--Panova
prove, in their partition labeling, at coordinate degree $7$,
\[
 \operatorname{mult}_{(34,6,2)}\mathbb C[\mathrm{Ch}_3^6]_7
 =7<8=
 \operatorname{mult}_{(34,6,2)}\mathbb C[\mathrm{Pow}_{3,4}^6]_7.
\]
Their Main Theorem (2a) and Section 3 additionally prove that every
partition occurring in the ambient coordinate ring occurs on
$\mathrm{Ch}_3^6$, in every coordinate degree. Thus no occurrence
obstruction separates $\mathrm{Pow}_{3,k}^6$ from this Chow variety,
for any $k$, although the displayed multiplicities separate them for
$k=4$ \cite{DIP2019Interface}. For the quantifiers we use the proof's
equation labeled \texttt{eq:posplethimpliesposchow}, which retains
the same partition and degree on both sides; the theorem's final
sentence has inconsistent indices. This is a source theorem at fixed
$(m,n)=(3,6)$, not a statement in all dimensions or about permanent
orbit closures. Our finite plethysm checker does not certify the
paper's Chow ranks or its all-degree no-occurrence proof.

Our concrete family has the stronger quotient $\rho_n$ already proved,
so neither an occurrence obstruction nor a multiplicity obstruction
in the direction $Z_n\not\subseteq\Omega_{4n+1}$ can arise.
This conclusion concerns precisely the coefficient-recoverable family
$\widetilde H_n$: setting $z=1$ and taking selector coefficients still
recovers every original cubic, while forgetting the selectors or
replacing the tuple by one product has the information losses proved
earlier. The displayed maps provide an actual endpoint in GCT orbit
geometry for the retained-fibre construction. They do not identify this
family with padded permanents or supply a Boolean satisfiability reduction.
